\scanpage{30}

\uline{\mspage{\uline{III} \unsure{0}}}

\uline{Introduction}

Dans le par. \unsure{4}, nous avons vu \deleted{que} lorsque
$X$ est propre sur $Y$ \deleted{loc. noeth, \unsure{I-adique}
\unsure{complet}, \inserted{\unsure{pour} $A$}}
\deleted{et si on \unsure{avait}} et si $Y'$ est une partie
fermée de $Y$, $X'=f^{-1}(Y')$, la \inserted{en particulier,}
commutation des foncteurs \deleted{$R^if$ et $R^i\hat f$ aux}
``images directes supérieures'' aux foncteurs complétion
[le long de $Y'$, resp $X'$].
Lorsque $Y$ était le spectre d'un anneau \deleted{\unsure{local}}
noeth $I$-adique complet, et $Y'=V(\unsure{I})$, ces résultats,
complétés par un th. d'existence (par 5), s'expriment
\deleted{en disant} \inserted{\unsure{en termes}}
\deleted{que $F'\unsure{\mathbin{-}|\longrightarrow}F'_{Y'}$ est un}
par l'équivalence de la catégorie des \deleted{faisceaux}
\inserted{Modules} coh. sur $\unsure{X'|X}$ et de la catégorie
des \unsure{pro-Modules} coh. sur $\unsure{X'_{Y'}|\widehat X_{Y'}}$,
\uline{compatible} avec les foncteurs des modules de cohomologie.

Dans ce par., nous \deleted{\unsure{généraliserons}}
\inserted{\unsure{voulons|montrons}}, que ces résultats \unsure{puissent}
se généraliser à certaines situations, où $f:X'\longrightarrow X$
n'est plus supposé propre [p. ex. $f$ une imm. ouverte\ldots];
la \inserted{seule} \unsure{différence} \unsure{essentielle}
est que \unsure{ici} \unsure{on se borne}
\deleted{\unsure{en ce qui concerne} des \unsure{énoncés}
à préciser \unsure{les} \unsure{th.} \unsure{de}
\unsure{connexité} et leurs \illegible{2}}
\inserted{\unsure{des bornes} \unsure{cohom.} qui
\unsure{pour des raisons} \unsure{seront}
$\unsure{\leq n}$, par \unsure{récurrence}.}
\unsure{du} \unsure{théor.} (cf N° 2).
La méthode sera celle des par 4 et 5,
\inserted{\unsure{d'ailleurs très générale}}
en utilisant des \uline{th. de finitude}
\inserted{(N° 1)} \deleted{\unsure{pour} les \unsure{cohomologies}
$\unsure{f_*}$ \unsure{qui se trouvent}}
\unsure{convenables}, faisant intervenir des propriétés
\deleted{\unsure{cohomologiques} \unsure{supposées}}
\inserted{s'expriment :}
locales des faisceaux envisagés, qui s'expriment :

\margin{29}
